\textbf{Participants and Experimental Task.} Thirty-nine participants completed 150 discrete choice tasks simulating a daily commute. On each trial, participants chose between three transport alternatives described by four attributes: cost, travel time, comfort, and CO\textsubscript{2} emissions. Attribute levels were fully combined and optimized using the ChoiceDesign Python package to minimize dominated alternatives and ensure an active processing baseline. To induce and test spatial habituation, the visual layout of the attributes remained perfectly stable for the first 100 trials. At trial 101, the spatial location of the attributes was unexpectedly disrupted, though the underlying economic options remained identical. The experiment was implemented in PsychoPy and conducted under controlled laboratory conditions.

\textbf{Eye-Tracking and Econometric Analysis.} Visual attention was recorded using Tobii Pro Spark eye-trackers sampling at 60 Hz. Fixation durations were computed by summing gaze events within 19 predefined Regions of Interest (ROIs), allowing us to calculate trial-specific gaze entropy and attribute-specific dwell times. After excluding trials with invalid timestamps or recording errors, the final dataset comprised 5,848 valid choice observations. Choice data and visual entropy were jointly estimated using a structural Latent Attention discrete choice model in R \cite{rstudio2020}. This model treats choice inertia as an endogenous function of visual search, using gaze entropy as a structural measurement equation to resolve observational equivalence. Full details regarding stimulus generation, eye-tracking calibration, ROI definitions, and the structural equation derivations are provided in the Supporting Information (SI) Appendix.
